\section{Joint real specialization and finite parameter data}\label{sec:real}
This section applies to a compact semialgebraic model $X$, a continuous
semialgebraic probability map $F$, and a continuous semialgebraic monic
polynomial $A_x$ of fixed degree $n$. The coefficients may be complex;
real and imaginary parts are always separate real coordinates. Fix a datum
$P\in F(X)$ and assume first that $A_x=A_0$ on $F^{-1}(P)$. Write
$A_0(z)=\prod_r(z-r)^{m_r}$. Choose disjoint isolating discs for these
roots. Compactness implies that all sufficiently small observation balls
have exactly $m_r$ target roots in the disc at $r$. Their monic factors
satisfy
\begin{equation}\label{eq:clusters}
 A_{r,x}(r+z)=z^{m_r}+\sum_{a<m_r}c_{r,a}(x)z^a,
 \qquad k_{r,a}=m_r-a.
\end{equation}
Polynomial multiplication and the isolating conditions define the unique
cluster factors semialgebraically. The differential of multiplication is
invertible for coprime factors, so the factors are analytic in the total
polynomial coefficients on this neighbourhood.

\subsection{Input representation and real-accessible covers}
\begin{definition}[Effective input]\label{def:input}
Real algebraic input consists of finite first-order polynomial formulas for
$X$ and the graphs of $F,A$, together with $P$, the design weights, and
cluster-isolating data. Every constant is represented by an integer
polynomial and a rational isolating interval selecting a real root.
Arithmetic in several such constants can equivalently be performed after
passing to a primitive element. A complex constant is a pair of represented
real algebraic numbers. Endpoints, clocks, polynomial coefficients, datum
probabilities, disc centres and squared radii are included in this rule.
All sign conditions, denominator exclusions and branch selections are
part of the formulas, not information supplied to a numerical oracle.
For a parameter family these are finite formulas in additional real
parameter coordinates; the numerical output assertion applies to algebraic
parameter values. No complexity bound is imposed.
\end{definition}

Adjoin the graph variables
\[
 R_{je}=\sqrt{w_j}(\sqrt{F(x)_{je}}-\sqrt{P_{je}}),\quad
 G=\sum_{j,e}R_{je}^2,\quad
 H_{r,a}=(\Re c_{r,a})^2+(\Im c_{r,a})^2.
\]
The nonnegative square-root choices are equations and inequalities on a
compact real graph $\Gamma$. Thus $G,H$ are polynomial coordinate
functions on $\Gamma$, even at zero observed cells.

\begin{definition}[Monomial and exact pieces]\label{def:monomial}
Real monomial data near $G=0$ are finitely many relatively compact
coordinate boxes on smooth real sources mapping properly, collectively,
onto a neighbourhood of $G=0$ in $\Gamma$. Each source component is of
one of the following two kinds. On an \emph{exact piece}, $G\equiv0$;
then all $H_{r,a}\equiv0$, and the piece is recorded separately without
entering any quotient of orders. On every other piece,
\begin{equation}\label{eq:monomial}
 G=u\prod_i|y_i|^{a_i},\qquad
 H_{r,a}\equiv0\quad\hbox{or}\quad
 H_{r,a}=v_{r,a}\prod_i|y_i|^{b_{r,a,i}}.
\end{equation}
The nonzero units are positive and bounded above and below on smaller
closed boxes. All exponents are nonnegative integers. A divisor is
recorded only if $a_i>0$ and it has a real point in the source box away
from the other divisors, with an admissible transversal on which $G>0$.
Identically zero $H$ contributes $+\infty$; $G\equiv0$ never contributes
$0/0$. Lower-dimensional exact and nonexact pieces are both retained.
\end{definition}

\begin{proposition}[A finite real cover]\label{prop:real-cover}
The graph $\Gamma$ has data as in Definition~\ref{def:monomial}.
The data cover every sufficiently small admissible observation
neighbourhood, and each recorded divisor supplies a real admissible arc.
For the input of Definition~\ref{def:input}, all pieces, integer orders,
and accessibility tests admit a finite effective construction.
\end{proposition}
\begin{proof}
We give the interfaces of the construction. First decompose the compact
semialgebraic graph into finitely many basic closed sets, using real
elimination \cite{BCR98,BPR06}. On each piece replace an inequality
$p(x)\ge0$ by $p(x)=s^2$. The original graph coordinates range over a
compact set and each slack is bounded by the maximum of $|p|^{1/2}$.
The resulting real algebraic lift is compact and its real projection is
exactly the original piece. This equality is the reason for using square
slacks before, rather than inequalities after, resolution.

Resolve each reduced algebraic lift in characteristic zero
\cite{BM97}. Over its nonsingular locus the resolution is an isomorphism,
so it covers every real point there. A real point not covered in this way
lies in the algebraic singular locus. Resolve that locus separately and
repeat. The dimension strictly decreases in each such recursion; after
at most the original dimension plus one levels every real point is
covered. This extra recursion is essential, for example at an isolated
real singular point whose complex branches have no real generic point.
Each map used is proper. The finite disjoint union of the resolved
sources is therefore a proper cover of the real lift.

On each smooth source separate identically zero functions first, including
$G$. Principalize the product of all remaining nonzero $G,H$ by blow-ups
with smooth centres. Since the factors are regular and their product has
normal crossings, each factor is a monomial times a unit in the same
coordinates. These blow-ups do not lose real points: the fibre above a
real point of a smooth centre is a real projective space, hence nonempty.
All source points still map into the square-slack lift and hence into
$\Gamma$. Take finitely many coordinate neighbourhoods of the compact
inverse image of $G=0$ and shrink them to smaller boxes. Properness shows
that their images contain $G\le\varepsilon_0$ in $\Gamma$ for some
$\varepsilon_0>0$: otherwise a sequence outside the images with $G\to0$
would have a lifted accumulation point above $G=0$. The units have the
stated bounds on these smaller boxes.

For each proposed divisor, test for a real point satisfying its equation,
the strict nonvanishing of the other divisor coordinates and units, and
the chart conditions. A transversal $y_i=s>0$ with the other coordinates
fixed then maps into $\Gamma$ and has $G>0$. If necessary reverse the
coordinate sign. Charts with no such point are omitted; their
lower-dimensional intersections have already been included in the cover.
All tests are first-order real formulas. Characteristic-zero constructive
resolution supplies finite algebraic equations for the centres and
charts; real quantifier elimination decides the tests and selects smaller
algebraic coordinate neighbourhoods. The recursion terminates by dimension,
and each resolution and principalization terminates by its classical
invariant. This proves both finite termination and coverage. The same
argument using real uniformization \cite{BM88} proves the existence
statement without choosing an effective presentation.
\end{proof}

\begin{corollary}[The scalar order data]\label{cor:orders}
For each nonrigid cluster coefficient define
\begin{equation}\label{eq:exponent}
 \beta_{r,a}=\min_{\text{recorded real divisors}}\frac{b_{r,a,i}}{a_i},
 \qquad \alpha=\min_{r,a}\frac{\beta_{r,a}}{k_{r,a}}.
\end{equation}
Either all coefficients vanish on an observation neighbourhood, or the
finite minimum $\alpha$ is positive rational. In the latter case
$\omega_P(t)\asymp t^\alpha$. Every finite coefficient minimum is
attained by an admissible transversal. The minima are independent of
the chosen real cover.
\end{corollary}
\begin{proof}
For a fixed $H$, let $\beta$ be the minimum in \eqref{eq:exponent}.
On each nonexact box, $b_i\ge\beta a_i$ for $a_i>0$, and $b_i\ge0$
otherwise. Bounded units give $H\le C G^\beta$. A finite cover makes
this uniform, so $|c_{r,a}|\le C h_w^\beta$. On a minimizing real
transversal, $G\asymp s^{a_i}$ and $H\asymp s^{b_i}$; inversion of
the positive analytic branch gives the reverse coefficient order.
Since $H=0$ on $G=0$, every relevant nonzero $H$ has $b_i>0$.
If no nonzero $H$ meets an observation divisor, the proper cover gives
an observation neighbourhood on which it vanishes.

For a monic polynomial, its maximal root modulus is comparable to
$\max_{a<m}|c_a|^{1/(m-a)}$. One direction follows from the elementary
symmetric functions; the other follows by summing a geometric series
outside a sufficiently large multiple of that maximum. The centre
belongs to every observation ball, so the diameter lies between the
maximal displacement and twice that displacement. This proves the root
order. Finally $\beta$ is the optimal exponent in a coefficient-versus-
observation inequality, an intrinsic property of the real graph. Two
covers therefore give the same number. This is the usual divisorial
presentation of a real \L{}ojasiewicz order, not a bound on ordinary
Taylor jets in the original coordinates.
\end{proof}

\subsection{The joint fibre and its exact diameter}
Assume the nonrigid case and write $\alpha=p/q$ in lowest positive terms.
For $\varepsilon>0$ let $\Psi_P(\varepsilon,x,U,V)$ assert membership in
$\Gamma$ together with
\begin{equation}\label{eq:weighted}
 R_{je}=\varepsilon^q U_{je},\qquad
 c_{r,a}=\varepsilon^{p k_{r,a}}V_{r,a},\qquad
 \sum_{j,e}U_{je}^2\le1.
\end{equation}
For complex coefficients the second equality is two real equations.
The leading joint set is specified without suppressing a quantifier:
\begin{align}
 (U,V)\in\cI_0(P)\quad\Longleftrightarrow\quad
 &\forall\rho>0\ \exists\varepsilon,x,U',V':\nonumber\\[-2pt]
 &0<\varepsilon<\rho,\quad
 \Psi_P(\varepsilon,x,U',V'),\quad
 \|(U',V')-(U,V)\|^2<\rho^2.                 \label{eq:closure-formula}
\end{align}
Here the centre $P$ is held fixed. Put $\cC_0(P)=\pr_V\cI_0(P)$.
Bounded model coordinates may be projected after closure: a convergent
subsequence of them exists, so this projection loses no limiting point.
The formula retains both the original inequalities and all joint
coefficient relations. Algebraic saturation by $\varepsilon$ is not a
substitute for \eqref{eq:closure-formula}.

\begin{theorem}[Joint specialization and root diameter]\label{thm:joint}
The set $\cC_0(P)$ is compact and is the Hausdorff limit of the rescaled
joint coefficient images of the radius-$t$ observation balls, with
$t=\varepsilon^q$. Define
\[
 Z_r(V)=Z\left(z^{m_r}+\sum_{a<m_r}V_{r,a}z^a\right).
\]
Then
\begin{equation}\label{eq:joint-diameter}
 \omega_P(t)=C_Pt^\alpha+o(t^\alpha),\qquad
 C_P=\max_{V,V'\in\cC_0(P)}\max_r d_\infty(Z_r(V),Z_r(V')),
 \quad 0<C_P<\infty.
\end{equation}
For real algebraic input there is a finite algorithm outputting the
rational exponent, quantifier-free formulas for $\cI_0,\cC_0$, and an
integer polynomial and rational isolating interval for $C_P$.
\end{theorem}
\begin{proof}
Corollary~\ref{cor:orders} bounds every $V$ in \eqref{eq:weighted};
$U$ is already bounded. For fixed positive $\varepsilon$ the projected
$V$-section is exactly the rescaled coefficient image. Its outer limit
is the closed compact set in \eqref{eq:closure-formula}. For a point
of that limit, distance to the positive section is a bounded
semialgebraic function of $\varepsilon$ with a zero subsequential
limit. One-variable semialgebraic monotonicity gives limit zero.
A finite net of the limit set proves the inner Hausdorff inclusion;
boundedness and compactness prove the outer inclusion. This proves
Hausdorff convergence, not merely the existence of individual arcs.

Continuity of monic root multisets on bounded coefficient sets now
gives the limiting root set. To justify the metric, let $\Delta$ be the
least separation of distinct base roots. For small $t$ each cluster
lies in its $\Delta/4$ disc. There is a within-cluster matching with
error tending to zero; any cross-cluster matching has an edge of length
at least $\Delta/2$. Thus optimal matchings cannot cross clusters.
Inside cluster $r$, translation by $r$ and dilation by $t^\alpha$
are exact in \eqref{eq:weighted}. The diameter of the joint root set
therefore has exactly the right-hand side of \eqref{eq:joint-diameter}
as its limit. Its positivity follows from an attaining coefficient
transversal and the root bound used in Corollary~\ref{cor:orders}.

For completeness we specify the last algorithmic interface. Quantifier
elimination applied to \eqref{eq:closure-formula} produces $\cI_0$ and
then its $V$-projection. Introduce labelled root coordinates through all
elementary symmetric identities of each monic cluster polynomial.
For two such root arrays $z,z'$ and $d\ge0$, the condition
\begin{equation}\label{eq:matching-formula}
 M_d(z,z'):=\bigvee_{\sigma\in\mathfrak S_m}
          \bigwedge_{i=1}^m |z_i-z'_{\sigma(i)}|^2\le d^2
\end{equation}
is exactly $d_\infty\le d$. Real and imaginary parts make these
polynomial inequalities. Exact equality to $d$ is $M_d$ together with
$\forall e\,(0\le e<d\Rightarrow\neg M_e)$. The maximum over the
finitely many clusters is encoded by conjunction for an upper bound
and by an attained equality in at least one cluster. Finally $C\ge0$
is the diameter if every pair of coefficient/root arrays has distance
at most $C$ and some pair has distance exactly $C$. Compactness gives
attainment. This is a first-order formula defining one real number over
the real algebraic input field. Elimination followed by real-root
isolation gives a polynomial over that field and an isolating interval;
a norm to $\Q$ and further isolation give the stated integer polynomial.
Ties between permutations or optimizing pairs do not require genericity.
No observation-ball maximum is an input to the order construction.
\end{proof}

\subsection{Parameter dependence}
A fixed-format family means a semialgebraic parameter set $\Theta$, a
semialgebraic $X\subset\Theta\times K$ with $K$ compact and each $X_\theta$
compact, continuous semialgebraic graphs $F_\theta,A_\theta$ on $X$, and a
semialgebraic choice $P_\theta\in F_\theta(X_\theta)$. The degree and the
number of observation coordinates are fixed. The graphs, including all
inequalities, have finite real algebraic presentations. Statements on a
stratum are uniform only on its compact subsets unless otherwise stated.

\begin{lemma}[Relative finite order lists]\label{lem:relative}
On the locus where the exact target fibre is a singleton, parameter space
has a finite semialgebraic refinement with fixed cluster multiplicities
and a fixed finite list of integer monomial orders and real-accessibility
flags as in Definition~\ref{def:monomial}. The source charts and units may
vary. Neither the divisors nor their labels are asserted to be canonical.
\end{lemma}
\begin{proof}
First apply semialgebraic triviality to the compact graph fibres. On
each resulting parameter cell the graph projection is proper: its
trivialization is a product with a compact fibre. Thus the graph is
relatively closed over that cell. Next partition by the degree,
multiplicities and an isolating choice for the base clusters. The distinct roots and the selected factors are Nash
functions after a finite semialgebraic refinement; alternatively include
their real and imaginary coordinates in a finite algebraic cover and
select its branches on cells. Real semialgebraic choice supplies positive
cluster-separation and observation radii; use smaller closed observation
neighbourhoods so that no root meets an isolating boundary. The lifted
cluster graph then has compact fibres. Apply elimination to this family
graph and its square-slack lifts.

We detail why resolution can be made a finite family procedure. On an
irreducible algebraic closure of a parameter cell, pass to its function
field. Resolve every dominating component of the generic lifted fibre,
then its successively singular subvarieties as in
Proposition~\ref{prop:real-cover}, and principalize the nonzero functions.
This uses finitely many equations and blow-ups. Spread these equations
out after clearing denominators. Delete a proper algebraic subset of the
base so that the centres are smooth relative to the base, the source
components and their relative singular loci have their generic
dimensions, and every displayed divisor has relative normal crossings.
These conditions are algebraic rank and nonvanishing conditions on
finitely many charts. Their bad-fibre images are constructible by
elimination and omit the generic point. Replacing those images by
their proper algebraic closures gives effective exceptional base loci. The maps remain proper because the blow-ups are
projective. The generic regular-locus isomorphisms and the decompositions
into regular and singular loci also spread out after this deletion.
Thus the recursive coverage proof applies to each remaining fibre, not
merely to the total space. Identically zero functions on a generic
component are recorded before spreading; a nonzero one stays nonzero
on the retained relative components after further shrinking.

Components not dominating this open base, and all excluded parameters,
lie over a proper algebraic subset. Repeat the construction there.
Noetherian induction on base dimension terminates. If a component or
cluster label requires a finite algebraic base cover, perform the
construction on that cover and then use real semialgebraic branch
selection; the finite numerical order lists descend without requiring
canonical labels. Finally the existence of a real generic divisor
point is a first-order condition on the parameter and chart variables.
Partition by these finitely many conditions. This last step discards
complex-only divisors and retains lower-dimensional real specializations.
It is not obtained by resolving the total space once and assuming that
resolution commutes with every specialization. The resulting finite
refinement has exactly the asserted lists and flags.
\end{proof}

\begin{theorem}[A finite real spectral atlas]\label{thm:atlas}
A fixed-format family admits a finite semialgebraic stratification with
the following properties.

\smallskip\noindent
\textup{(i)} The exact target fibres have constant semialgebraic
topological type and dimension on each stratum. On a nonidentified
stratum, $\omega_\theta(t)\to\diam Z(A(F_\theta^{-1}(P_\theta)))>0$.

\smallskip\noindent
\textup{(ii)} On the identified locus, each stratum is either spectrally
rigid on an observation neighbourhood or has a fixed positive rational
exponent given by \eqref{eq:exponent}. Its cluster multiplicities,
accessible integer order list and minimizing index set are fixed.

\smallskip\noindent
\textup{(iii)} On every nonrigid identified stratum, the joint coefficient
leading sets and their root images are semialgebraic families with
constant semialgebraic topological type after refinement. Their positive
diameters $C(\theta)$ are Nash functions after refinement, and
\[
 \omega_\theta(t)=C(\theta)t^\alpha+O_H(t^{\alpha+\nu})
\]
on each compact $H$ of a sufficiently refined stratum, for some $\nu>0$.
The exponent $\nu$ need not agree between strata.

\smallskip\noindent
\textup{(iv)} All these objects have finite real-algebraic descriptions.
On a subfamily for which a finite Fisher-cone normal form is proved,
the stratification can additionally retain its active supports and
leading-constant choices. Outside such a subfamily the constant program
is semialgebraic, not asserted to be quadratic.
\end{theorem}
\begin{proof}
The graph of the exact target fibre is a semialgebraic projection with
compact fibres. Semialgebraic triviality \cite{Hardt80} gives (i)'s
partition; marking its singleton fibres is a first-order condition.
For a fixed centre, decreasing compact observation balls converge to
the exact fibre in Hausdorff distance by subsequence compactness. This
proves the asserted nonidentified limit.

On the singleton locus apply Lemma~\ref{lem:relative}. The finitely many
integer lists and fixed multiplicities give (ii), including rigidity
and fixed minimizing indices. On a stratum with $\alpha=p/q$, use
\eqref{eq:closure-formula} with $\theta$ held fixed throughout its
quantifiers. Quantifier elimination gives semialgebraic families
$\cI_0(\theta),\cC_0(\theta)$. In particular one must not take closure
while letting $\theta$ vary: that different construction would insert
limits from neighbouring strata. Apply triviality to these sets and
their root graphs. The compact matching formula gives a semialgebraic
positive function $C$, which is Nash on a finite stratification.

We explain the uniform remainder assertion. For each fixed parameter,
Theorem~\ref{thm:joint} gives Hausdorff convergence of bounded rescaled
root images to the leading root image. Its error is a semialgebraic
function of $(\theta,t)$. A finite cylindrical decomposition of its
graph gives finitely many polynomial relations for its nonzero
one-variable branches. The possible positive leading orders are among
the finitely many slopes determined by pairs of monomials in these
relations; vanishing coefficient specializations are separated by
further sign conditions. Choose $\nu$ smaller than the least positive
order in each piece. The existence of $M<\infty$ and $t_0>0$ with
error at most $Mt^\nu$ for $0<t<t_0$ is first-order and true pointwise.
Semialgebraic choice supplies such $M(\theta),t_0(\theta)$; refine until
they are continuous. On compact subsets they have the required upper
and positive lower bounds. The same refinement bounds rescaled root
sets and separates base clusters. Diameters differ by at most twice
Hausdorff distance, proving (iii). Finally a finite list of rational
Fisher/KKT sign conditions can be added to any of these strata, proving
(iv). This refinement connects the valuation and Fisher partitions
without identifying two different kinds of walls.
\end{proof}

\begin{remark}
The scalar part of Theorem~\ref{thm:atlas} uses classical resolution and
semialgebraic geometry. Its formulation records the joint target data
needed by a spectral inverse problem. It is not a claim that exact-fibre
homeomorphism type determines the exponent, that the exponent determines
the leading set, or that the local leading set detects a remote fibre.
Sections~\ref{sec:multiplicity} and \ref{sec:wall} compute these distinctions
rather than suppressing them.
\end{remark}
